\chapter{绪论}\label{chap:introduction}

操作系统连接硬件资源与应用程序，其实现牵涉特权级、地址空间、异常中断、进程调度、文件接口和设备驱动等多个层面。对于计算机科学与技术专业的本科教学而言，学生不仅需要理解进程、页表、系统调用等概念，还需要在一套可以运行、可以修改的代码中看到这些概念如何发生联系。没有这样的实例，很多机制会停留在术语层面。代码规模适中、结构清楚、便于调试的教学操作系统，正是弥合这一距离的有效材料。

本课题以同济大学 Unix V6++ 教学操作系统为基础，研究其面向 RISC-V 平台的移植与 Rust 重构。重构前的 Unix V6++ 面向 IA-32 架构，内核主体采用 C++ 编写；本文以保留 Unix V6 风格系统调用、进程语义和 Shell 使用方式为前提，重新实现启动、中断异常、系统调用、进程管理、设备适配和用户程序支持等路径。文件系统和内存管理作为系统运行的支撑条件出现，不作为展开重点。最终目标是让系统能够通过 QEMU virt 机器启动到可交互 Shell，并运行常见用户程序、文件命令、进程和信号测试。

\section{研究背景与意义}

操作系统课程的难点在于其研究对象并非单个算法或单个程序，而是由硬件机制和软件抽象共同组成的整体系统。进程调度依赖时钟中断，系统调用依赖特权级切换，Shell 命令依赖用户程序装载和文件描述符，终端输入又会牵动设备中断、TTY 和信号处理。各部分之间存在紧密联系。缺少能够运行的内核实例，学生很难理解这些机制之间的相互作用。

经典教材通常从概念层面对操作系统进行系统论述，例如进程管理、内存管理、文件系统、输入输出和安全机制等内容\cite{tanenbaum2017modern,arpaci2019ostep}。这些教材为理解操作系统原理提供了必要框架，但对于本科阶段的系统实践而言，还需要一个规模可控、便于阅读和修改的内核作为实验对象。Unix V6 及其教学化衍生系统正是基于这一需求而长期被采用：它们功能相对完整，又不像现代通用操作系统那样庞大复杂，适合作为理解内核内部结构的切入点。

同济大学 Unix V6++ 是面向操作系统教学开发的 Unix V6 风格实验系统。该系统将早期 Unix 的核心思想迁移到更容易获得的硬件平台，并使用面向对象方式重构了进程、文件系统、设备和中断等模块。其教学价值在于保留了 Unix 系统“麻雀虽小，五脏俱全”的特点：学生可以通过阅读和修改系统代码，观察系统调用如何进入内核、进程如何被调度、文件如何通过 inode 和缓冲区访问磁盘，以及终端输入如何驱动 Shell 交互。陈闳中教授关于 Unix V6++ 教学实践的回顾也强调了“让学生真正动手修改、调试内核”的课程目标\cite{chenhz_unixv6pp}。

近年，RISC-V 开源指令集架构在教学、科研与产业中被持续采用。RISC-V 架构具有开放、简洁、模块化和可扩展等特点，其特权架构规范对异常、中断、页表和特权级切换等操作系统所需机制进行了明确描述\cite{riscv_privileged_20211203}。与传统封闭指令集相比，RISC-V 更适合作为系统软件教学中的硬件抽象基础：学生既可以学习指令集和特权级机制，也可以进一步理解操作系统如何利用这些机制实现地址隔离、异常处理和设备访问。

基于上述背景，将 Unix V6++ 移植到 RISC-V 平台具有教学和工程两重意义。课程层面，它把原有 Unix V6++ 实验习惯与 RISC-V 开放生态放到同一套系统中；工程层面，移植工作要求重新审视启动流程、地址布局、中断控制器、设备模型、用户程序 ABI 和构建系统。学生面对的不再是孤立知识点，而是一条从固件入口到 Shell 命令返回的执行路径。

\section{国内外研究现状}

教学操作系统的设计目标通常不是追求工业级完备性，而是在有限代码规模内覆盖操作系统的核心机制。Unix V6 是这一方向的重要源头。它以简洁的代码实现了进程、文件、目录、设备和系统调用等基本机制，长期被用作理解 Unix 操作系统结构的经典材料。以 Unix V6 为思想来源的教学系统强调可读性和可修改性，适合用来展示操作系统各模块之间的配合关系。

MIT 的 xv6 是当前国际上应用较广的 Unix-like 教学操作系统。xv6 早期版本面向 x86，后续版本迁移到 RISC-V，并在 MIT 操作系统课程中用于讲解进程、页表、锁、文件系统和系统调用等主题\cite{xv6_book_2023}。xv6 代码量较小，结构清晰，注释和实验材料也较完备，许多高校将其作为操作系统课程的重要参考。不过，xv6 更强调从 Unix V6 思想出发重写一个小型教学内核，而 Unix V6++ 则更突出在本校课程体系中长期积累的教学接口、实验习惯和用户态环境。

除 xv6 外，OS/161、Pintos 等系统也被广泛用于操作系统教学。OS/161 常用于配合同步、虚拟内存和文件系统实验，Pintos 则强调从半成品内核出发逐步完成调度、用户程序和文件系统等功能。这类系统能够训练学生的内核编程能力，但其设计背景、运行环境和实验组织方式各不相同。对于本课题而言，更重要的是依托已有 Unix V6++ 教学基础完成平台迁移，而不是重新选择一个全新的教学内核。

国内 RISC-V 操作系统教学也已经积累了较多实践。rCore 教程以 Rust 语言和 RISC-V 架构为基础，从批处理系统、地址空间、进程管理到文件系统逐步构造操作系统内核\cite{rcore_tutorial_v3}。该项目说明了 Rust 与 RISC-V 结合用于系统软件教学的可行性，也展示了借助类型系统和模块化组织改善内核代码可维护性的思路。与此同时，QEMU 为 RISC-V 提供了稳定的系统级仿真环境，能够模拟 virt 机器、串口、PLIC、VirtIO 设备等硬件组件，降低了教学实验对真实硬件的依赖\cite{qemu_riscv_doc}。

Rust 语言近年来也逐渐进入系统软件开发领域。Rust 通过所有权、生命周期和类型检查等机制，不依赖垃圾回收即可提供内存安全能力\cite{rust_reference}。内核开发仍然需要处理裸机启动、外设访问、页表修改和上下文切换等底层操作，无法完全避免 unsafe 代码，但 Rust 的模块系统、枚举、模式匹配和 trait 等语言机制有助于表达复杂状态和抽象边界。对于教学内核而言，Rust 既能保留接近硬件的控制能力，也能减少空指针、悬垂引用和资源释放错误带来的理解负担。

这些教学系统共同说明了小型内核对操作系统课程的价值，RISC-V 与 Rust 也提供了新的实现基础。问题在于，Unix V6++ 并不能通过简单替换编译目标完成迁移：启动入口、页表、中断、设备、用户程序格式和构建链都会随平台改变。本文以保留原有教学语义为前提，处理这些具体迁移问题。

\section{本文主要研究内容}

本文围绕 Unix V6++ 的 RISC-V 移植与 Rust 重构展开，主要研究内容包括以下四项。

第一，完成系统启动和地址空间基础的迁移。重构后的内核运行于 QEMU virt 平台，通过 OpenSBI 进入 Rust 内核入口，完成 BSS 段清零、串口初始化、页表建立和内核初始化。内存管理部分围绕后续进程运行所需的 Sv39 页表、内核高地址映射、用户空间映射和 MMIO 映射展开，并接入页帧分配器，为内核对象、用户进程图像和内核栈提供页级资源。

第二，完成中断、异常和系统调用框架的重构。RISC-V 平台不再使用 IA-32 的 GDT、IDT、8259A 等机制，而是依赖 \texttt{stvec}、\texttt{scause}、\texttt{sepc}、\texttt{sstatus} 等特权寄存器完成 Trap 处理。本文实现了统一 Trap 入口，在汇编层保存通用寄存器和关键控制状态，并在 Rust 层分发系统调用、时钟中断、外部中断和异常。系统调用接口保留 Unix V6 风格编号和语义，由用户态 C 库封装 \texttt{ecall} 指令进入内核，内核再根据系统调用号分发到文件系统、进程管理和终端等模块。

第三，完成进程管理和用户程序支持。重构后的系统实现了进程控制块、进程状态机、独立内核栈、任务上下文切换、sleep/wakeup 和调度逻辑。用户程序采用 ELF64/RISC-V 格式，内核在 \texttt{exec} 过程中解析 ELF 头和程序头，将 text、data 和 stack 区域映射到用户空间，并构造 argc/argv 用户栈。系统实现了 \texttt{fork}、\texttt{exec}、\texttt{wait}、\texttt{exit} 等核心进程系统调用，同时支持基本信号机制，使终端中断、非法访问和用户发送信号能够作用于运行中的进程。

第四，完成设备、终端和用户交互环境的适配。文件系统继续提供 Unix V6++ 的 inode、目录、文件描述符和缓冲区接口，本文主要关注其与进程、系统调用和设备层的连接方式。设备层由原先面向 PC 平台的 ATA/IDE 和键盘显示路径，调整为 QEMU virt 平台上的 UART、PLIC 和 VirtIO Block。终端部分基于串口输入输出和 TTY 行规程支撑 Shell，进而验证命令执行、管道、重定向和常见用户程序。

与重构前的 main 分支相比，当前系统的主要变化体现在三个层面：体系结构从 IA-32 转向 RISC-V，内核主体从 C++ 转向 Rust，运行环境从传统 PC 设备模型转向 QEMU virt、OpenSBI、PLIC、UART 和 VirtIO 组成的仿真平台。第 7 章的构建与测试验证表明，这些变化已经贯通到用户库、Shell、文件系统镜像和 QEMU 启动流程，而不只停留在内核局部模块。

\section{本文组织结构}

本文共分为八章。第 1 章为绪论，介绍课题背景、研究意义、相关教学操作系统发展现状以及本文主要研究内容。第 2 章分析原 Unix V6++ 系统结构与重构目标，说明当前系统相对 main 分支的总体变化。第 3 章介绍 RISC-V 启动流程、地址空间设计、Sv39 页表和物理内存管理实现。第 4 章介绍 Trap 框架、中断异常处理和系统调用机制。第 5 章介绍进程模型、调度、核心进程系统调用、ELF 用户程序加载和信号机制。第 6 章介绍 Unix V6 文件系统适配、VirtIO 块设备、缓冲区管理、TTY 和 Shell 交互环境。第 7 章介绍系统构建、QEMU 运行环境和功能测试结果。第 8 章总结全文工作，分析当前实现的不足，并展望后续可改进方向。
